% =====================================================================
% cheatsheet.tex — KACTL-style multi-column competitive-programming
% cheatsheet template.
%
% This file is a TEMPLATE consumed by tools/cheatsheet.py. The script
% replaces three placeholder tokens (each written as two percent signs,
% a NAME, then two percent signs) before compiling:
%
%   TITLE       -> the cheatsheet title (LaTeX-escaped)
%   NUMCOLUMNS  -> number of multicol columns (integer)
%   BODY        -> the generated body (one block per algorithm)
%
% The tokens themselves appear only in the title/multicols area below.
% It compiles with a plain `pdflatex` and only common packages:
% article + geometry + multicol + listings + xcolor.
% =====================================================================
\documentclass[8pt,a4paper,landscape]{extarticle}

\usepackage[utf8]{inputenc}
\usepackage[T1]{fontenc}
\usepackage[a4paper,landscape,margin=1cm]{geometry}
\usepackage{multicol}
\usepackage{xcolor}
\usepackage{listings}

% Tighten everything up for a dense reference sheet.
\setlength{\parindent}{0pt}
\setlength{\columnsep}{1.2em}
\setlength{\columnseprule}{0.2pt}
\pagestyle{empty}

% --- Colours -----------------------------------------------------------
\definecolor{cs@rule}{gray}{0.55}
\definecolor{cs@kw}{RGB}{0,0,153}
\definecolor{cs@cmt}{RGB}{0,128,0}
\definecolor{cs@str}{RGB}{163,21,21}
\definecolor{cs@bg}{gray}{0.97}

% --- Listings: small, tight, monospaced --------------------------------
\lstset{
  basicstyle=\ttfamily\scriptsize,
  keywordstyle=\color{cs@kw}\bfseries,
  commentstyle=\color{cs@cmt}\itshape,
  stringstyle=\color{cs@str},
  showstringspaces=false,
  breaklines=true,
  breakatwhitespace=false,
  columns=fullflexible,
  keepspaces=true,
  tabsize=2,
  aboveskip=2pt,
  belowskip=2pt,
  lineskip=-1pt,
  xleftmargin=0pt,
  frame=none,
  backgroundcolor=\color{cs@bg},
  % `listings` is byte-oriented, so multi-byte UTF-8 chars that show up in
  % comments must be mapped explicitly or inputenc rejects them.
  extendedchars=true,
  literate=
    {á}{{\'a}}1 {é}{{\'e}}1 {í}{{\'i}}1 {ó}{{\'o}}1 {ú}{{\'u}}1
    {Á}{{\'A}}1 {É}{{\'E}}1 {Í}{{\'I}}1 {Ó}{{\'O}}1 {Ú}{{\'U}}1
    {ñ}{{\~n}}1 {Ñ}{{\~N}}1 {ü}{{\"u}}1 {Ü}{{\"U}}1 {ç}{{\c c}}1
    {—}{{---}}1 {–}{{--}}1 {…}{{\ldots}}1
    {¿}{{?`}}1 {¡}{{!`}}1 {«}{{<<}}1 {»}{{>>}}1
    {→}{{$\rightarrow$}}1 {←}{{$\leftarrow$}}1
    {≤}{{$\leq$}}1 {≥}{{$\geq$}}1 {≠}{{$\neq$}}1
    {×}{{$\times$}}1 {·}{{$\cdot$}}1 {∞}{{$\infty$}}1 {°}{{$^\circ$}}1,
}

% Section heading for one algorithm.
\newcommand{\csalgo}[1]{%
  \vspace{0.4ex}%
  {\bfseries\small #1}\par
  \vspace{-0.6ex}{\color{cs@rule}\hrule height 0.4pt}\vspace{0.4ex}%
}
% Stats line (complexity / use case) under an algorithm heading.
\newcommand{\csstats}[1]{{\scriptsize\itshape #1\par\vspace{0.3ex}}}

\begin{document}

% --- Title -------------------------------------------------------------
\begin{center}
  {\Large\bfseries Chuletario ICPC CUNEF}
\end{center}
\vspace{0.5ex}
{\color{cs@rule}\hrule height 0.6pt}
\vspace{0.8ex}

\begin{multicols}{3}
\csalgo{Binary search}
\csstats{time: O(log n), space: O(1)  |  Encontrar un valor en un array ordenado.}
\begin{lstlisting}[language=C++]
#include <vector>
using namespace std;
int binary_search_idx(const vector<int>& a, int target) {
    int lo = 0, hi = (int)a.size() - 1;
    while (lo <= hi) {
        int mid = lo + (hi - lo) / 2;
        if (a[mid] == target) return mid;
        if (a[mid] < target) lo = mid + 1;
        else hi = mid - 1;
    }
    return -1;
}
\end{lstlisting}

\csalgo{Fast I/O (C++)}
\begin{lstlisting}[language=C++]
// Fast I/O — put these two lines at the very start of main():
ios::sync_with_stdio(false);
cin.tie(nullptr);
// Then read/write as usual (now unbuffered from C stdio, much faster):
//   int x; cin >> x;          // input
//   cout << x << "\n";         // output ('\n' is faster than endl)
\end{lstlisting}

\csalgo{Fast I/O (Python)}
\begin{lstlisting}[language=Python]
import sys

data = sys.stdin.buffer.read().split()   # fast input: read all tokens at once
it = iter(data)
def read_int():                          # next integer from the input
    return int(next(it))

out = sys.stdout.write                   # fast output: out(str(x) + "\n")
\end{lstlisting}

\csalgo{Dijkstra}
\csstats{time: O((V+E) log V), space: O(V+E)  |  Camino más corto desde un origen con pesos no negativos.}
\begin{lstlisting}[language=C++]
#include <vector>
#include <queue>
#include <functional>
using namespace std;
const long long INF = 1e18;
vector<long long> dijkstra(int n, const vector<vector<pair<int, int>>>& adj, int src) {
    vector<long long> dist(n, INF);
    priority_queue<pair<long long, int>, vector<pair<long long, int>>, greater<>> pq;
    dist[src] = 0;
    pq.push({0, src});
    while (!pq.empty()) {
        auto [d, u] = pq.top();
        pq.pop();
        if (d > dist[u]) continue;
        for (auto [v, w] : adj[u]) {
            if (dist[u] + w < dist[v]) {
                dist[v] = dist[u] + w;
                pq.push({dist[v], v});
            }
        }
    }
    return dist;
}
\end{lstlisting}

\csalgo{Fenwick Tree (BIT)}
\csstats{build: O(n), update: O(log n), query: O(log n), space: O(n)  |  Sumas de prefijos con actualizaciones puntuales.}
\begin{lstlisting}[language=C++]
#include <vector>
using namespace std;
struct FenwickTree {
    int n;
    vector<int> tree;
    FenwickTree(int n) : n(n), tree(n) {}
    void update(int index, int delta) {
        for (index++; index <= n; index += index & -index)
            tree[index - 1] += delta;
    }
    int query(int index) {
        int acc = 0;
        for (; index > 0; index -= index & -index)
            acc += tree[index - 1];
        return acc;
    }
    int query_range(int l, int r) { return query(r) - query(l); }
};
\end{lstlisting}

\csalgo{Convex hull}
\csstats{time: O(n log n), space: O(n)  |  El menor polígono convexo que contiene un conjunto de puntos.}
\begin{lstlisting}[language=C++]
#include <vector>
#include <algorithm>
using namespace std;
struct P {
    long long x, y;
    bool operator<(const P& o) const {
        return x < o.x || (x == o.x && y < o.y);
    }
};
long long cross(const P& O, const P& A, const P& B) {
    return (A.x - O.x) * (B.y - O.y) - (A.y - O.y) * (B.x - O.x);
}
vector<P> convex_hull(vector<P> pts) {
    int n = pts.size();
    if (n < 3) return pts;
    sort(pts.begin(), pts.end());
    vector<P> h(2 * n);
    int k = 0;
    for (int i = 0; i < n; i++) {
        while (k >= 2 && cross(h[k - 2], h[k - 1], pts[i]) <= 0) k--;
        h[k++] = pts[i];
    }
    for (int i = n - 2, t = k + 1; i >= 0; i--) {
        while (k >= t && cross(h[k - 2], h[k - 1], pts[i]) <= 0) k--;
        h[k++] = pts[i];
    }
    h.resize(k - 1);
    return h;
}
\end{lstlisting}
\end{multicols}

\end{document}
